% Part: first-order-logic
% Chapter: lindstrom
% Section: lindstrom-proof

\documentclass[../../../include/open-logic-section]{subfiles}

\begin{document}

\olfileid{mod}{lin}{prf}

\olsection{لنڈسٹروم (Lindstr\"om) دا قضیہ}

\begin{lem}
\ollabel{lem:lindstrom}
فرض کرو تجریدی منطق معمول مطابق اے، تے $!E \in L(\Lang{L})$، جتھے $\Lang{L}$ محدود
اے، تے ایہ وی فرض کرو کہ کوئی $n \in \Nat$ موجود اے
ایس طرح کہ ہر دو !!{structure}s $\Struct{M}$ تے~$\Struct{N}$
لئی جے $\Struct{M} \equiv_n
\Struct{N}$ تے $\Struct{M} \models_L !E$ ہووے، تاں
$\Struct{N}
\models_L !E$ وی ہووے۔ تاں $!E$ پہلے درجے دے اک
!!{sentence} دے ہم ارز اے؛ یعنی پہلے درجے دا کوئی $!D$
موجود اے جس لئی $\Mod[L](L){!E} = \Mod[L](L){!D}$؛
ایہ دونوں نمونہ طبقے دتی زبان دیاں ساختاں لئی نیں۔
\end{lem}

\begin{proof} 
$n$ اوہ ہووے جس لئی ہر دو $n$-ہم ارز !!{structure}s
$\Struct{M}$ تے $\Struct{N}$، $!E$ نوں دتی قدر اُتے
متفق ہون۔ \olref[bas][pis]{prop:qr-finite} یاد کرو:
محدود !!{language} وچ مقدار نمائی درجہ~$n$ توں
ودھ نہ ہون والے پہلے درجے دے !!{sentence}s، منطقی
ہم ارزی تک، صرف محدود تعداد وچ نیں۔ ایہ لفظی جملیاں دی محدود
تعداد نہیں: ہر منطقی ہم ارزی جماعت وچوں اک نمائندہ چنو۔ معمول
مطابق منطق وچ، پچھلی وضاحت مطابق، عام پہلے درجے دے جملے اوہناں
دے معمول دے صدق نال شامل سمجھے جا سکدے نیں۔ ہن ہر مقرر
!!{structure}~$\Struct{M}$ لئی $!D_{\Struct{M}}$،
چنے پہلے درجے دے !!{sentence}s~$!A$ دا اوہ عطف
ہووے جو $\Struct{M}$ وچ سچ ہون تے جنہاں لئی
$\QuantRank{!A} \le n$؛ ایہ عطف ہن واقعی محدود اے۔ ایس طرح
$\Struct{N} \models !D_{\Struct{M}}$ جے تے صرف جے
$\Struct{N}
\equiv_n \Struct{M}$۔ فیر $!D = \textstyle\bigvee
\Setabs{!D_{\Struct{M}}}{\Struct{M} \models_L !E}$
رکھو؛ ایس فصل وچ وی ہر حاصل درجہ وار قسم لئی صرف اک
نمائندہ رکھو، سو ایہ واقعی محدود پہلے درجے دا جملہ اے۔ خالی
عطف نوں ہمیشہ سچا جملہ تے خالی فصل نوں ہمیشہ جھوٹا جملہ سمجھو۔

نتیجہ $\Mod[L](L){!E} = \Mod[L](L){!D}$ نکلدا اے۔
واقعی، جے $\Struct{N} \models_L !D$، تاں کسے
$\Struct{M} \models_L
!E$ لئی $\Struct{N} \models !D_{\Struct{M}}$ ہووے گا،
جس توں $\Struct{N} \models_L !E$ وی نکلدا اے
(لمے دے مفروضے نال)۔ الٹے پاسے، جے
$\Struct{N} \models_L !E$، تاں $!D_\Struct{N}$
$!D$ دا اک جزوِ انفصال اے، تے چونکہ
$\Struct{N}
\models !D_\Struct{N}$، ایس لئی
$\Struct{N} \models_L !D$ وی۔
\end{proof}

\begin{thm}[Lindstr\"om دا قضیہ]
  \ollabel{thm:lindstrom} فرض کرو $\tuple{L, \models_L}$ وچ
  معمول مطابق منطق دیاں شرطاں، محدود گواہی تے نزولی
  L\"owenheim--Skolem دیاں خاصیتاں نیں۔ تاں
  $\tuple{L, \models_L} \le \tuple{F, \models}$ (اس لئی
  $\tuple{L, \models_L}$ پہلے درجے دی منطق دے ہم ارز اے)۔
\end{thm}

\begin{proof}
ہر $!E \in L(\Lang{L})$ لوو۔ توسیع دی خاصیت توں اک
محدود !!{language} $\Lang{L_0}\subseteq\Lang{L}$ ملدی اے
جس دے گھٹائے روپ اُتے $!E$ دا صدق منحصر اے۔ ایہ صدقی وابستگی
اے؛ اسیں $!E$ نوں چھوٹی زبان دا جملہ نہیں منّاں گے۔ ہر چھوٹی
زبان دی ساخت نوں اصل زبان تک ودھایا جا سکدا اے، تے اکّو گھٹایا
روپ رکھن والیاں ایس طرح دیاں ساختاں وچ $!E$ دا صدق اکّو اے۔
سو \olref{lem:lindstrom} دے محدود نمائندیاں والے استدلال نوں
ایس صدقی خاصیت اُتے ورتنا کافی اے: کوئی $n\in\Nat$ ایسا ہووے
کہ دو اصل !!{structure}s $\Struct{M}$ تے~$\Struct{N}$ دے
چھوٹے زبان والے روپ $\Struct{M}\equiv_n\Struct{N}$ ہون
تاں $!E$ بارے متفق ہون۔ ایتھے ایہ ہم ارزی صرف انہاں چھوٹے
روپاں لئی لکھی گئی اے۔ محدود درجہ وار قسم دے اوہو عطف تے فصل
اک پہلے درجے دا !!{sentence} بنان گے جو چھوٹی زبان وچ اے
تے اصل زبان وچ وی $!E$ دے ہم ارز اے۔ ایس توں
$\tuple{L, \models_L} \le \tuple{F, \models}$ نکلے گا؛
معمول مطابق منطق وچ پہلے درجے دی شمولیت الٹی سمت تے قوساں وچ
لکھی ہم ارزی وی دیندی اے۔ ہِٹھاں محدود خالص تعلقی زبان لئی
\olref[bas][pis]{thm:b-n-f} مطابق مقدار نمائی درجے دی ہم ارزی نوں
اگے-پیچھے دی محدود الجبری شرط $I_n(\emptyseq,\emptyseq)$
نال ورتاں گے۔

تضاد لئی فرض کرو کہ ہر $n$ لئی اصل !!{structure}s
$\Struct{M}_n$ تے $\Struct{N}_n$ موجود نیں جنہاں دے
چھوٹے روپ مقدار نمائی درجہ $n+1$ تک ہم ارز نیں، پر
$\Struct{M}_n\models_L !E$ تے
$\Struct{N}_n\not\models_L !E$۔ اک نواں مستقل $c_\circ$
چنو۔ پہلی ساخت وچ اک عنصر چن کے، محدود درجے دی ہم ارزی دی
اگے والی شرط نال دوسری وچ ایسا متناظر عنصر چن سکدے آں کہ
انہاں نوں نواں مستقل نام دین پچھوں محدود زبان
$\Lang{L_0}\cup\{c_\circ\}$ دے روپ مقدار نمائی درجہ $n$ تک ہم ارز
ہون۔ اصل زبان دے باقی مستقلاں نوں ایس نامی عنصر دی قدر دیو،
تے باقی محمول خالی رکھو۔ چھوٹا روپ نہیں بدلیا، سو $!E$ دا
مختلف صدق وی نہیں بدلیا۔ ہن اسی محدود ودھی زبان لئی
$I_n(\emptyseq,\emptyseq)$ حاصل اے۔ اوہدیاں بند جوہری
!!{sentence}s صرف محدود تعداد وچ نیں؛ اس لئی کسے بے حد
ذیلی سلسلے وچ سارے جوڑیاں دے اوہو جوہری !!{sentence}s سچ ہون گے۔ جوڑی
دے دونوں پاسے درجہ صفر اُتے وی متفق نیں۔ ایہ ذیلی سلسلہ
ایویں دوبارہ گنو کہ نئی جوڑی نمبر $n$ دا پرانا درجہ کم از کم
$n$ ہووے؛ اصل $I_n$ شرط قائم رہندی اے۔ ہم شکل نقلاں وچ
مستقلاں دیاں متناظر قدراں نوں مشترک نامی حصہ بنا کے پہچانو،
تے ہر جوڑی دے باقی عناصر نوں الگ نشان دیو۔ بند جوہری صدق دا
اتفاق مشترک نامی حصے اُتے برابری تے محمول تعبیر دا اتفاق یقینی
کردا اے؛ سو اتحاد وچ اصل محمول دوہاں پاسیاں نال رل سکدے نیں۔
$\Struct{M}$ اوہ !!{structure} ہووے جو ساریاں $\Struct{M}_n$ ساختاں دا اتحاد اے،
جس وچ انہاں دیاں اصل محمول تعبیراں دا اتحاد اے تے ہور مخلوط
ٹپلاں اُتے محمول جھوٹے نیں؛ ہر $\Struct{M}_n$ اصل زبان
دے روپ وچ ذیلی ساخت اے۔ \olref[lsp]{thm:abstract-p-isom}
دے درست کیتے ثبوت وانگ، $\Struct{M}^*$ دی
!!{domain} $\Domain{M}\cup\Domain{M}^{<\omega}$ وچ
لڑیّاں دے الگ نشان والے کوڈ ورتو۔ ودھی !!{language} وچ
جوڑن دے محمول $P$ تے~$Q$ اوہو نیں۔

اسی طرح $\Struct{N}_n$، $\Struct{N}$ تے
$\Struct{N}^*$ مقرر کرو۔ بڑی !!{structure} نوں
$\Struct{C}$ آکھو، تاکہ اصل پاسے والی ساخت دا ناں نہ رلے۔
اوہدی !!{domain} وچ $\Struct{M}^*$ تے $\Struct{N}^*$
دے !!{domain}s، فطری عدد~$\Nat$ دے وکھرے کوڈ تے اوہناں
دی فطری ترتیب~$\le$، نال جزوی نقشیاں دے الگ کوڈ شامل ہون۔
یک مقامی نواں محمول $Z$ عددی کوڈ چنّے۔ اصل نامی حصہ مشترک
اے؛ لڑیّاں، عددی کوڈ تے نقشہ کوڈ اصل عناصر توں وکھرے نیں۔
اِس دی !!{language} وچ وادھو محمول !!{domain}s
$\Domain{M}$، $\Domain{N}$، $\Domain{M}^{<\omega}$ تے
$\Domain{N}^{<\omega}$ نوں چنّن، تے $R$ تے $S$،
$\Struct{M}_n$ تے $\Struct{N}_n$ دے دائرے تے متعلقہ لڑی
کوڈ ایس معنی وچ چنّن:
\begin{align*}
  \Domain{M_n} & = \Setabs{a \in \Domain{M}}{R(a, n)}; & 
  \Domain{N_n} & = \Setabs{a \in \Domain{N}}{S(a,n)}; \\
  \Domain{M}^{<\omega}_n & = \Setabs{a \in \Domain{M}^{<\omega}}{R(a,n)}; &
  \Domain{N}^{<\omega}_n & = \Setabs{a \in \Domain{N}^{<\omega}}{S(a,n)}. 
\end{align*}
ایس !!{structure} وچ ہن اگے-پیچھے دی باقی گہرائی تے اصل جوڑی دا اشاریہ وکھ رکھنا
لازم اے۔ اصل عنصر دائرے چنّن والے محمول $U_0$ تے $V_0$
آکھو، تے دو مقامی محمول ایویں مقرر کرو:
$U(a,m)\leftrightarrow(U_0(a)\land R(a,m))$ تے
$V(b,m)\leftrightarrow(V_0(b)\land S(b,m))$۔ ہن لڑی
کوڈ انہاں چنے حصیاں وچ نہیں آندے۔ نواں تین مقامی رشتہ $J(d,m,h)$ دا مطلب ہووے:
کوڈ $h$ والا نقشہ جوڑی نمبر $m$ وچ گھٹ توں گھٹ $d$
ہور قدم اگے-پیچھے توسیع برداشت کردا اے۔ تین مقامی محمول
$G(h,a,b)$ اوس نقشے دا گراف اے۔ اصل $I_d$ وچ آنے والے
محدود tuple-نقشیاں نال مشترک نامی حصے دا شناختی نقشہ ملا کے
ایسے گراف بناؤ، جیویں پچھلے ثبوت وچ کیتا۔ محدود درجے دی شرط توں
ہر $m$ لئی $J(m,m,h)$ والا کوئی کوڈ ملدا اے؛ گہرائی صفر
تک ہر توسیع اصل جوڑی $m$ نوں قائم رکھدی اے۔ ایہ کوڈ بندی
اصل لڑیّاں دی تعمیر نوں نہیں بدل دی، پر غلط اکّو اشاریے والے
رشتے دی تھاں اصل جزوی نقشیاں نوں صریح رکھدی اے۔

ہن !!{language}~$\Lang{L}$ نوں $R$، $S$، $J$، $G$،
$Z$ وغیرہ نال ودھاؤ؛ تمام کوڈ بندی دیاں علامتاں اصل زبان
توں نئی چنّیاں نیں۔ نسبتی روپ، بولین بندش تے مسوّر دی
خاصیت نال اک تجریدی !!a{sentence}~$!D$ حاصل اے جو ہر
عدد کوڈ $m$ لئی $\Struct{M}_m\models_L !E$ تے
$\Struct{N}_m\not\models_L !E$ آکھدا اے۔ ایس نوں
بنان لئی عارضی نواں مستقل جوڑی دا اشاریہ بنے؛ پہلے $U$
تے $V$ دے اوس پیرامیٹر والے حصیاں تک $!E$ تے اوہدی نفی
دا نسبتی روپ لوو، فیر «جے اشاریہ عددی اے» دی شرط نال ساریاں
قدراں اُتے عمومی مسوّر لوو۔ ایہ پیرامیٹر ذیلی ساخت اندر ہون
ضروری نہیں؛ ہر پرانے مستقل دی قدر اندر ہونا ضروری اے۔ نال
پہلے درجے دے جملیاں دا سیٹ $\Delta$ ورتو۔ ایہ $Z$ اُتے
$\le$ نوں خطی ترتیب بناندا اے، اک پہلا عددی رکن چنّدا اے،
ہر عدد نوں فوری جانشین تے ہر غیر پہلے عدد نوں فوری پیش رو
دیندا اے۔ $B(e,d)$ فوری پیش رو دا رشتہ ہووے، جس دی تعریف
سیٹ وچ ایہ جملہ کرے:
\[
\forall e,d\,(B(e,d)\leftrightarrow
(Z(e)\land Z(d)\land e<d\land
\neg\exists t\,(Z(t)\land e<t\land t<d))).
\]
خطی ترتیب اُتے فوری پیش رو تے جانشین ایس تعریف نال اکیلے
مقرر ہوندے نیں۔ ہر عددی جوڑی
لئی $U$ تے $V$ دے اصل-عنصر حصے غیر خالی ہون، تے ہر
پرانے مستقل تے $c_\circ$ دی قدر دوہاں حصیاں وچ ہووے۔ ہر
$J(d,m,h)$ والا گراف انہاں حصیاں وچ جزوی یک بہ یک نقشہ
ہووے، ساریاں نامی قدراں نوں اپنی قدر اُتے ای ملائے، تے ہر
اصل محمول دے ٹپلاں دا صدق دونوں رخ محفوظ رکھے۔ ایہ گراف
شرطاں پچھلے ثبوت والی پہلے درجے دی شرطاں دا جوڑی-اشاریے
والا روپ نیں۔ صریح طور اُتے ایہ جملے شامل کرو:
\[
\begin{gathered}
\forall m\,(Z(m)\to(\exists a\,U(a,m)\land\exists b\,V(b,m))),\\
\forall d,m,h\,(J(d,m,h)\to(Z(d)\land Z(m))),\\
\forall d,m,h,a,b\,((J(d,m,h)\land G(h,a,b))\to(U(a,m)\land V(b,m))),\\
\forall h,a,b,b'\,((G(h,a,b)\land G(h,a,b'))\to b=b'),\\
\forall h,a,a',b\,((G(h,a,b)\land G(h,a',b))\to a=a').
\end{gathered}
\]
ہر پرانے مستقل $c$، تے $c_\circ$، لئی ایہ جملے شامل کرو:
\[
\forall m\,(Z(m)\to(U(c,m)\land V(c,m))),\qquad
\forall d,m,h\,(J(d,m,h)\to G(h,c,c)).
\]
ہر پرانے $r$-مقامی محمول $A$ لئی ایہ جملہ شامل کرو:
\[
\forall d,m,h,a_1,\ldots,a_r,b_1,\ldots,b_r\,
\bigl((J(d,m,h)\land\bigwedge_{i=1}^{r}G(h,a_i,b_i))\to
(A(a_1,\ldots,a_r)\leftrightarrow A(b_1,\ldots,b_r))\bigr).
\]
لڑی کوڈ حصیاں نوں اصل-عنصر حصیاں نال نہ ملاؤ۔ سیٹ وچ ایہ
ابتدائی کوڈ تے دونوں توسیع شرطاں وی شامل ہون:
\[
\begin{gathered}
\forall m\,(Z(m)\to\exists h\,J(m,m,h)),\\
\forall d,m,h,e,a\,((J(d,m,h)\land B(e,d)\land U(a,m))\to
\exists k,b\,(J(e,m,k)\land V(b,m)\land G(k,a,b)\land
\forall u,v\,(G(h,u,v)\to G(k,u,v)))),\\
\forall d,m,h,e,b\,((J(d,m,h)\land B(e,d)\land V(b,m))\to
\exists k,a\,(J(e,m,k)\land U(a,m)\land G(k,a,b)\land
\forall u,v\,(G(h,u,v)\to G(k,u,v)))).
\end{gathered}
\]
ایتھے $U$ تے $V$ صرف اصل عناصر چنّن، تے $J$ دی شرط وچ
پہلے دو دلائل عددی ہون۔ $\Delta$ وچ ہر غیر متعلقہ پرانا محمول خالی تے ہر
غیر متعلقہ پرانا مستقل $c_\circ$ دے برابر ہون دے جملے وی
شامل کرو۔ ایہ سارے جملے تے $!D$، $\Struct{C}$ وچ سچ
نیں۔ نال عددی مستقل $\nu_k$، معیاری عدد $k\in\Nat$ دے
کوڈ نوں نام دیندے ہون؛ $\Delta$ وچ پہلے رکن توں انہاں دی
فوری جانشینی لڑی دے جملے رکھو۔ اسیں صرف پہلے درجے دا سیٹ
آکھ رہے آں، اک واحد محدود جملہ یا حقیقی معیاری عددی دائرہ
پہلے درجے نال مجبور کرن دا دعویٰ نہیں۔

اک نواں مستقل $c_*$ چن کے جملہ-سیٹ
$\{!D\}\cup\Delta\cup\{Z(c_*),\ \nu_k<c_*:k\in\Nat\}$
لوو، جتھے سخت ترتیب $<$، $\le$ تے نا برابری نال تعریف
ہوندی اے۔ ہر محدود ذیلی سیٹ $\Struct{C}$ وچ پورا ہوندا
اے: $c_*$ نوں اوہدے اندر آئے سارے محدود معیاری عددی کوڈاں
توں وڈا معیاری عدد دیو؛ باقی سیٹ پہلے ای سچ اے۔ زبان ودھاون
نال $!D$ دا صدق وی قائم اے۔ محدود گواہی دی خاصیت توں
ایس پورے سیٹ دا نمونہ $\Struct{C}^{\dagger}$ ملدا اے۔
اوہدے اندر $c_*$ دی قدر $n^*$ سارے معیاری عددی کوڈاں توں
وڈی اے۔ اصل زبان والے، غیر خالی ذیلی !!{structure}s
$\Struct{M_{n^*}}$ تے $\Struct{N_{n^*}}$، متعلقہ
اصل-عنصر حصیاں توں لوو؛ $!D$ توں
$\Struct{M_{n^*}}\models_L !E$ تے
$\Struct{N_{n^*}}\not\models_L !E$۔ ہر معیاری محدود
$k$ لئی $n^*-k$ دا مطلب فوری پیش رو دے $k$ قدم اے؛
ایہ موجود اے تے $n^*-k>0$، کیوں جے پہلی قدر تک محدود
قدم وچ پہنچنا $n^*$ نوں کسے معیاری کوڈ دے برابر بنا دیندا۔
ہن $\mathcal{I}$ اوہناں محدود نقشیاں $p$ دا سیٹ ہووے جنہاں
دا دائرہ $k$ عناصر دا اے تے جو کسے گراف $G(h,-,-)$ دی
پابندی نیں، جتھے $J(n^*-k,n^*,h)$ اے۔ اصل جوڑی دا
اشاریہ ہن ہمیشہ $n^*$ اے؛ صرف باقی گہرائی گھٹدی اے۔
ابتدائی کوڈ والی شرط خالی نقشہ دیندی اے، سو خاندان غیر خالی
اے۔ جے مطلوب عنصر پہلے ای دائرے یا حاصل قدراں وچ اے، اوہو
$p$ رکھو۔ نہیں تاں لکھی توسیع شرط نال پورا پرانا گراف قائم
رکھ کے مطلوب نواں جوڑا لوو، تے صرف $p$ تے اوہ جوڑا رکھو؛
ایہ $k+1$ عناصر والا محدود نقشہ اے تے اگلی پیش رو گہرائی
دا کوڈ اوہنوں رکھدا اے۔ گراف دیاں شرطاں یک بہ یک پن، محمول تے
مستقل حفاظت دیندیاں نیں۔ ایس طرح دونوں اگے-پیچھے رخ حاصل
نیں تے $\Struct{M_{n^*}}\simeq_p\Struct{N_{n^*}}$۔
پر \olref[lsp]{thm:abstract-p-isom}، جس دے معمول مطابق ہون تے
نزولی مفروضے ایتھے دونوں موجود نیں، انہاں نوں $L$ دے لحاظ
نال جملہ وار ہم ارز کردا اے؛ ایہ $!E$ دے مختلف صدق دے
خلاف اے۔ سو مقدار نمائی درجے دی محدود حد موجود اے، تے اوپر دا محدود
نمائندیاں والا استدلال ہر اصل زبان دے ہر جملے لئی مطلوب
پہلے درجے دا مساوی جملہ دیندا اے۔
\end{proof}

\end{document}
